% TOP_PROBLEM: {"id": 7200025, "title": "The covering problem of Rado", "category_id": 6, "category": {"id": 6, "name": "geometry", "display_name": "Geometry", "description": "Euclidean and non-Euclidean geometry, geometric structures.", "slug": "geometry", "order_index": 6, "created_at": "2026-07-31T15:26:25.671Z"}, "status": "open"}
% FIRST_POSTED: null
% ENTRY_KIND: statement_only
% Imported from: batch19.tex; UnsolvedMath ID 7200025
\documentclass[11pt]{article}
\usepackage{fontspec}
\setmainfont{Latin Modern Roman}
\setsansfont{Latin Modern Sans}
\usepackage{amsmath,amssymb,mathtools}
\usepackage{unicode-math}
\setmathfont{Latin Modern Math}
\usepackage{xurl}
\usepackage[hidelinks]{hyperref}
\newcommand{\sref}[2]{\hyperlink{source:#1:#2}{[S#2]}}
\newcommand{\cataloguescope}[1]{\par\noindent\textbf{Question.} #1\par}
\begin{document}
% UnsolvedMath ID 7200025; source catalogue code AMR-071-0025
\setcounter{section}{7200024}
\section{The covering problem of Rado}\label{7200025}
{\small\sffamily UnsolvedMath ID 7200025 · Geometry}
\subsection{Definitions and mathematical statement}

\paragraph{Shared notation.}$\mathbb N=\{1,2,\ldots\}$, and $\mathbb Z,\mathbb Q,\mathbb R,\mathbb C$ denote the integers, rationals, reals and complex numbers. Any other conventions are specified locally.


Let $\mathcal C$ be a nonempty finite collection of closed axis-parallel squares of positive side length in $\mathbb R^2$, with side lengths allowed to differ. Write $|E|$ for planar Lebesgue area. A subcollection $\mathcal S\subseteq\mathcal C$ is disjoint if $Q\cap Q'=\varnothing$ whenever $Q,Q'\in\mathcal S$ are distinct; boundary contact therefore counts as intersection. Define
\[M(\mathcal C)=\max_{\substack{\mathcal S\subseteq\mathcal C\\\mathcal S\text{ disjoint}}}\left|\bigcup_{Q\in\mathcal S}Q\right|,\qquad c_\square=\inf_{|\bigcup_{Q\in\mathcal C}Q|=1}M(\mathcal C).\]
\paragraph{Statement recorded in the supplied source.}
Determine $c_\square$: how small can the largest area of a disjoint subcollection be when the union of the original squares has area one? Equivalently, determine the largest universal $c$ for which every finite collection satisfies $M(\mathcal C)\geq c\,|\bigcup_{Q\in\mathcal C}Q|$. \sref{7200025}{2}
\subsection{Short English statement}
Among finite collections of axis-parallel squares covering area one, determine the smallest possible maximum area that can be retained by selecting pairwise disjoint squares.
\subsection{Sources}
\begin{description}
\item[\hypertarget{source:7200025:1}{S1}] ULAM.AI and the UnsolvedMath Contributors, UnsolvedMath: A Curated Collection of Open Mathematics Problems, record 7200025 (AMR-071-0025). CC BY 4.0. \url{https://huggingface.co/datasets/ulamai/UnsolvedMath}.
\item[\hypertarget{source:7200025:2}{S2}] Gian Maria Dall’Ara and Adrian Dumitrescu, Rado’s covering problem for cubes and balls: a semi-survey (2026). Section 2, equation (3), pp. 2–3. \url{https://arxiv.org/pdf/2604.17509v2}.
\end{description}
\label{end:7200025}
\end{document}
